\section{Introduzione}

Il pricing di strumenti derivati e' uno dei problemi classici della finanza quantitativa.
Nel caso delle opzioni europee, il modello di Black-Scholes fornisce una formula chiusa per il prezzo teorico sotto ipotesi precise sulla dinamica del sottostante, sul mercato e sulla volatilita' \cite{black1973pricing}.
Questa disponibilita' di una soluzione analitica rende Black-Scholes un ambiente ideale per valutare in modo controllato l'accuratezza di metodi numerici e di modelli di machine learning.

La domanda di progetto e':

\begin{quote}
Una rete neurale feed-forward riesce ad approssimare accuratamente la funzione di pricing Black-Scholes per opzioni call europee?
\end{quote}

Il progetto confronta tre approcci:
\begin{itemize}
    \item formula analitica di Black-Scholes, usata come ground truth;
    \item simulazione Monte Carlo, usata come benchmark numerico;
    \item rete neurale feed-forward, usata come modello supervisionato e pricing surrogate.
\end{itemize}

L'interesse principale non e' sostituire Black-Scholes nel caso in cui la formula chiusa sia disponibile, ma verificare se una rete neurale possa apprendere una funzione di pricing nota.
Questo e' un primo passo utile per comprendere l'uso di modelli neurali in contesti piu' complessi, dove formule analitiche non sono disponibili o risultano computazionalmente costose.
